\usepackage[T1]{fontenc}
\usepackage[utf8]{inputenc}
\usepackage{lmodern}
\usepackage{microtype}
\usepackage[margin=28mm,headheight=14pt]{geometry}
\usepackage{amsmath,amsthm,amssymb,mathtools}
\usepackage{booktabs,array,float}
\usepackage{xcolor}
\usepackage[font=small,labelfont=bf]{caption}
\usepackage{fancyhdr}
\usepackage{enumitem}
\usepackage[colorlinks=true,linkcolor=black,citecolor=black,urlcolor=black,hypertexnames=false]{hyperref}
\usepackage{cleveref}

\definecolor{accent}{RGB}{24,99,111}
\definecolor{accentlight}{RGB}{120,170,178}
\definecolor{pale}{RGB}{234,243,244}
\definecolor{gapred}{RGB}{151,57,49}
\definecolor{ink}{RGB}{51,51,51}
\definecolor{midgray}{RGB}{140,140,140}

\pagestyle{fancy}
\fancyhf{}
\fancyhead[L]{\small Christopher Yamba}
\fancyhead[R]{\small Four diagonals for doubled letters, five for tripled}
\fancyfoot[C]{\thepage}
\renewcommand{\headrulewidth}{0.3pt}
\setlength{\parindent}{1.2em}
\setlength{\parskip}{0pt plus 1pt}
\setlist[enumerate]{itemsep=3pt,topsep=5pt}
\setlist[itemize]{itemsep=2pt,topsep=4pt}
\newlist{parts}{enumerate}{1}
\setlist[parts]{label=\textup{(\roman*)},ref=\textup{(\roman*)},leftmargin=*,itemsep=2pt,topsep=4pt}
\numberwithin{equation}{section}

\theoremstyle{plain}
\newtheorem{mainthm}{Theorem}
\renewcommand*{\themainthm}{\Alph{mainthm}}
\providecommand*{\theHmainthm}{}
\renewcommand*{\theHmainthm}{\Alph{mainthm}}
\newenvironment{maintheorem}[1]{\setcounter{mainthm}{\numexpr`#1-65\relax}\begin{mainthm}}{\end{mainthm}}
\newcommand*{\mainthmstarletter}{}
\newtheorem*{mainthmstar}{Theorem~\mainthmstarletter}
\newenvironment{maintheorem*}[1]{\renewcommand*{\mainthmstarletter}{#1}\begin{mainthmstar}}{\end{mainthmstar}}
\newtheorem{theorem}{Theorem}[section]
\newtheorem{lemma}[theorem]{Lemma}
\newtheorem{proposition}[theorem]{Proposition}
\newtheorem{corollary}[theorem]{Corollary}
\newtheorem{conjecture}[theorem]{Conjecture}
\newtheorem{question}[theorem]{Question}
\theoremstyle{definition}
\newtheorem{definition}[theorem]{Definition}
\newtheorem{example}[theorem]{Example}
\theoremstyle{remark}
\newtheorem{remark}[theorem]{Remark}

\crefname{mainthm}{Theorem}{Theorems}
\crefname{theorem}{Theorem}{Theorems}
\crefname{lemma}{Lemma}{Lemmas}
\crefname{proposition}{Proposition}{Propositions}
\crefname{corollary}{Corollary}{Corollaries}
\crefname{conjecture}{Conjecture}{Conjectures}
\crefname{question}{Question}{Questions}
\crefname{definition}{Definition}{Definitions}
\crefname{example}{Example}{Examples}
\crefname{remark}{Remark}{Remarks}
\crefname{section}{Section}{Sections}
\crefname{appendix}{Appendix}{Appendices}
\crefname{figure}{Figure}{Figures}
\crefname{table}{Table}{Tables}
\crefname{equation}{equation}{equations}
\crefname{partsi}{part}{parts}
\newcommand{\thmpart}[2]{\cref{#1}\textup{(#2)}}

\newcommand{\NN}{\mathbb{N}}
\newcommand{\ZZ}{\mathbb{Z}}
\newcommand{\QQ}{\mathbb{Q}}
\newcommand{\RR}{\mathbb{R}}
\DeclarePairedDelimiter{\floor}{\lfloor}{\rfloor}
\DeclarePairedDelimiter{\ceil}{\lceil}{\rceil}
\DeclarePairedDelimiter{\abs}{\lvert}{\rvert}
\newcommand{\iver}[1]{[#1]}
\newcommand{\pos}[1]{\left(#1\right)_{+}}
\newcommand{\interval}[2]{\{#1,\ldots,#2\}}
\newcommand{\word}[1]{\ensuremath{\mathtt{#1}}}
\newcommand{\gbinom}[2]{\genfrac{[}{]}{0pt}{}{#1}{#2}_{q}}
\newcommand{\Ex}{\mathbb{E}}

\newcommand{\cont}{\mathbf{n}}
\newcommand{\contsize}{\lvert\mathbf{n}\rvert}
\newcommand{\Words}[1]{\mathcal{W}(#1)}
\newcommand{\upto}{\mathinner{\ldotp\ldotp}}
\newcommand{\cnt}[1]{\lvert w\rvert_{#1}}
\newcommand{\cntof}[2]{\lvert #1\rvert_{#2}}
\newcommand{\cntr}[2]{\lvert w\rvert_{#1\upto #2}}
\newcommand{\Pmat}[1]{\Psi(#1)}
\newcommand{\psitwo}{\psi_2}
\newcommand{\pairvec}{\mathbf{x}}
\newcommand{\fibre}[1]{\psi_2^{-1}(#1)}
\newcommand{\amb}[1]{d(#1)}
\newcommand{\dmax}[1]{d_{\max}(#1)}
\newcommand{\rev}[1]{#1^{\mathrm{R}}}
\newcommand{\comp}[1]{#1^{\mathrm{C}}}
\newcommand{\rc}[1]{#1^{\mathrm{RC}}}
\newcommand{\proj}[2]{A^{(#1)}(#2)}
\newcommand{\projset}[2]{\mathcal{A}_{#1}(#2)}
\newcommand{\Bmat}[1]{\mathbf{B}_{#1}}
\newcommand{\one}{\mathbf{1}}
\newcommand{\Zalt}{Z}

\newcommand{\qx}{q_x}
\newcommand{\qy}{q_y}
\newcommand{\zmax}{z_{\max}}
\newcommand{\zmin}{z_{\min}}
\newcommand{\vals}[1]{\mathcal{Z}(#1)}
\newcommand{\Aset}[1]{\mathcal{A}(#1)}
\newcommand{\Cset}[1]{\mathcal{C}(#1)}
\newcommand{\Abal}{A^{\mathrm{bal}}}
\newcommand{\Cbal}{C^{\mathrm{bal}}}
\newcommand{\Acon}{A^{\mathrm{con}}}
\newcommand{\Ccon}{C^{\mathrm{con}}}
\newcommand{\pre}[2]{#1_{\le #2}}
\newcommand{\suf}[2]{#1_{> #2}}
\newcommand{\gset}[2]{\gamma(#1,#2)}
\newcommand{\ruleE}{\textup{(E)}}
\newcommand{\ruleN}{\textup{(N)}}
\newcommand{\ruleCh}{\textup{(Ch)}}
\newcommand{\ruleG}{\textup{($\Gamma$)}}
\newcommand{\TTwo}{\textup{(T2)}}
\newcommand{\TFive}{\textup{(T5)}}
\newcommand{\TFiveR}{\textup{(T5$^{\mathrm{R}}$)}}
\newcommand{\Lam}[1]{\Lambda_{#1}}
\newcommand{\Ups}[1]{\Upsilon_{#1}}
\newcommand{\minor}[3]{M_{#1,#2,#3}}
\newcommand{\Rig}{\operatorname{Rig}}
\newcommand{\Tmat}[1]{\mathbf{T}_{#1}}

\newcommand{\sgp}[2]{\langle #1,#2\rangle}
\newcommand{\simplex}[1]{\mathcal{T}_{#1}}
\newcommand{\levset}[2]{\mathcal{L}_{#1}(#2)}
\newcommand{\levpoly}[2]{\Phi^{(#1)}_{#2}}
\newcommand{\Gram}{\mathbf{G}}
\newcommand{\cstar}{c^{*}}
\newcommand{\goldphi}{\varphi}
\newcommand{\Euler}[1]{E_{#1}}
\newcommand{\Fib}[1]{F_{#1}}
\newcommand{\Luc}[1]{L_{#1}}
\newcommand{\Kv}{K^{\mathrm{v}}}
\newcommand{\Kh}{K^{\mathrm{h}}}

\newcommand{\lrle}{\le_{\mathrm{lr}}}
\newcommand{\Kern}{\mathcal{K}}
\newcommand{\QP}{\operatorname{QP}}
\newcommand{\projP}{\mathsf{P}}
\newcommand{\projQ}{\mathsf{Q}}
\newcommand{\labw}{\mathsf{s}}
\newcommand{\rodd}{r^{\mathrm{odd}}_{\infty}}
\newcommand{\reven}{r^{\mathrm{even}}_{\infty}}

\newcommand{\checkfile}[1]{\textup{\texttt{checks/\detokenize{#1}}}}
\newcommand{\thesisp}[1]{printed p.~#1, PDF p.~\number\numexpr#1+4\relax}
\newcommand{\thesispp}[2]{printed pp.~#1--#2, PDF pp.~\number\numexpr#1+4\relax--\number\numexpr#2+4\relax}
